% from ega3-7-fr.tex
\oldpage[III]{210}
où, pour tout $\mathscr O_X$-Module cohérent $\mathscr F$, plat sur $Y$, on
désigne par $\operatorname{EP}(\mathscr F;y)$ (ou
$\operatorname{EP}(f,\mathscr F;y)$ ou $\operatorname{EP}(Y,\mathscr F;y)$) la
fonction $\operatorname{EP}(\mathscr Q_\bullet;y)$ correspondant au complexe
$\mathscr Q_\bullet$ dont le seul terme $\ne0$ est de degré $0$ et égal à
$\mathscr F$. On voit donc qu'on peut se ramener à étudier les caractéristiques
d'Euler-Poincaré de complexes réduits à un seul terme.
\end{env}

\begin{proposition}[7.9.7]
\label{III.7.9.7.fr}
Sous les hypothèses de (7.9.4), soient $Y'$ un préschéma localement noethérien,
$g:Y'\to Y$ un morphisme,
\[
  X'=X\times_Y Y',
  \qquad f'=f_{(Y')}:X'\longrightarrow Y',
\]
$\mathscr P'_\bullet$ le complexe fini
\[
  \mathscr P_\bullet\otimes_{\mathscr O_Y}\mathscr O_{Y'}
\]
de $\mathscr O_{X'}$-Modules ; $\mathscr P'_\bullet$ est formé de
$\mathscr O_{X'}$-Modules cohérents et $Y'$-plats, et pour tout $y'\in Y'$, on a
\[
  \operatorname{EP}(\mathscr P'_\bullet;y')
  =\operatorname{EP}(\mathscr P_\bullet;g(y')).
  \tag{7.9.7.1}\label{III.7.9.7.1.fr}
\]
\end{proposition}

\begin{proof}
Les $\mathscr O_{X'}$-Modules $\mathscr P'_i$ étant images réciproques des
$\mathscr P_i$ par la projection $X'\to X$ sont cohérents, ils sont $Y'$-plats en
vertu de ($0_{\mathrm I}$, 6.2.1) et (I, 4.14.5), la question étant locale sur
$X$, $Y$ et $Y'$ ; enfin, on sait que $f'$ est propre (II, 5.4.2), donc le
premier membre de (7.9.7.1) est défini. La formule (7.9.7.1) résulte alors de
(6.10.4.2), (7.7.2) et du lemme (7.6.7), en se ramenant, comme on peut toujours
le faire, au cas où $Y$ et $Y'$ sont affines.
\end{proof}

\begin{proposition}[7.9.8]
\label{III.7.9.8.fr}
Supposons vérifiées les hypothèses de (7.9.4) et en outre qu'il existe un entier
$i_0$ tel que
\[
  \mathscr T_i(\kappa(y))=0
  \qquad\text{pour $i\ne i_0$ et tout $y\in Y$.}
\]
Alors
\[
  \mathscr T_{i_0}(\mathscr O_Y)
  =\mathscr H^{-i_0}(f,\mathscr P_\bullet)
\]
est un $\mathscr O_Y$-Module localement libre, dont le rang en $y\in Y$ est égal
à $(-1)^{i_0}\operatorname{EP}(f,\mathscr P_\bullet;y)$.
\end{proposition}

\begin{proof}
Remarquons d'abord que les hypothèses de (7.4.4) sont vérifiées par les
$\mathscr T_j^{\mathscr O_y}$, donc (7.4.7) leur est applicable, et l'hypothèse
entraîne que $\mathscr T_i^{\mathscr O_y}$ est \emph{nul} pour $i\ne i_0$ en
vertu de (7.5.3) ; en raison de (7.3.3), $\mathscr T_{i_0}$ est donc aussi exact,
et par suite (7.8.4), $\mathscr H^{-i_0}(f,\mathscr P_\bullet)$ est localement
libre et son rang en un point $y\in Y$ est
\[
  \operatorname{rang}_{\kappa(y)}
  \mathscr T_{i_0}(\kappa(y))
  =\operatorname{EP}(f,\mathscr P_\bullet;y)
\]
par définition, puisque $\mathscr T_i(\kappa(y))=0$ pour $i\ne i_0$.
\end{proof}

\begin{corollary}[7.9.9]
\label{III.7.9.9.fr}
Soient $Y$ un préschéma localement noethérien, $f:X\to Y$ un morphisme propre,
$\mathscr F$ un $\mathscr O_X$-Module cohérent et $Y$-plat ; on suppose qu'il
existe un entier $i_0$ tel que
\[
  H^i\bigl(f^{-1}(y),
  \mathscr F\otimes_{\mathscr O_Y}\kappa(y)\bigr)=0
  \qquad\text{pour tout $i\ne i_0$ et tout $y\in Y$.}
\]
Alors $\mathrm R^{i_0}f_*(\mathscr F)$ est un $\mathscr O_Y$-Module localement
libre, dont le rang en $y$ est égal à
$(-1)^{i_0}\operatorname{EP}(f,\mathscr F;y)$.
\end{corollary}

En particulier :

\begin{corollary}[7.9.10]
\label{III.7.9.10.fr}
Sous les conditions préliminaires de (7.9.9) pour $X$, $Y$ et $\mathscr F$,
supposons que l'on ait $\mathrm R^if_*(\mathscr F)=0$ pour tout $i>0$. Alors
$f_*(\mathscr F)$ est un $\mathscr O_Y$-Module localement libre, dont le rang en
$y$ est égal à $\operatorname{EP}(f,\mathscr F;y)$.
\end{corollary}

\begin{proof}
Il suffira, en vertu de (7.9.9) de prouver le lemme suivant :
\end{proof}

\begin{lemma}[7.9.10.1]
\label{III.7.9.10.1.fr}
Sous les hypothèses de (7.9.10), on a
\[
  H^i\bigl(f^{-1}(y),
  \mathscr F\otimes_{\mathscr O_Y}\kappa(y)\bigr)=0
\]
pour tout $i>0$ et tout $y\in Y$.
\end{lemma}

\begin{proof}
En effet, on peut se borner au cas où $Y=\operatorname{Spec}(A)$ est affine.
Avec les notations de (7.9.4), et $\mathscr P_\bullet$ étant réduit à son terme
de degré $0$ égal à $\mathscr F$, on a en effet
$\mathscr T_p(\mathscr O_Y)=0$ pour $p<0$ par hypothèse ; on conclut de (7.3.7)
que $\mathscr T_p$ est exact pour $p<0$, et le lemme résulte alors de
l'équivalence de (7.7.5, a)) et (7.7.5, d)).
\end{proof}

